% ============================================================
% Conclusion and Future Directions
%
% PURPOSE:
%   Summarize your contributions, revisit your research questions,
%   and propose directions for future work. Do NOT introduce
%   new results or arguments not already in the thesis body.
%
% TIPS:
%   - Be direct: "This thesis showed X" not "This thesis explored X."
%   - Map each contribution back to a research question.
%   - Future work should be concrete and specific, not just
%     "more research is needed."
%   - This is often the last thing a committee member reads —
%     make it clear and confident.
%
% TARGET LENGTH: 3--5 pages.
% ============================================================

\chapter{Conclusion and Future Directions}
\label{chapter_conclusion}

\noindent
[Write a brief opening paragraph restating what this thesis set
out to do and what it accomplished.]


\section{Summary of Contributions}

\noindent
[List your main contributions clearly. Be specific — say what you
demonstrated, built, or established, not just what you explored.]

\begin{enumerate}
    \item{} [State your first contribution concisely.]

    \item{} [State your second contribution concisely.]

    % Add or remove items as needed.
\end{enumerate}


\section{Answers to Research Questions}
% If your thesis used objectives rather than questions,
% rename this section "Achievement of Objectives."

\noindent
\textbf{RQ1: [Restate your first research question.]}
\newline

[State what your thesis found in answer to this question.
1--3 sentences.]
\newline

\noindent
\textbf{RQ2: [Restate your second research question.]}
\newline

[State what your thesis found in answer to this question.]


\section{Future Work}

\noindent
[Describe specific, actionable directions for future research.
Connect these to the limitations of your current work. "Future
work could extend this study to larger datasets" is more useful
than "further research is needed."]
